\subsection{I-27. Filtering the actual Hoza masks does not restore a short coherent code}
\label{I-27}

\paragraph{The precise repair being tested.}
I-23 required every seed of a code to preserve the Grassmann cone.
One could instead keep only those seeds that happen to preserve it.
For the actual base-six weighting, this filtering has an exact
description and cannot combine noncoherent bits into a coherent
action.  This is a statement about that formula and formal
cocharacters; it is not a lower bound for arbitrary nonlinear
decoders or dense interpolation maps.

Let $2\le r\le w-2$, let
$C:\binom{[w]}r\longrightarrow\{0,1\}^{\ell}$ be any binary
code, and put
\[
 \rho_s(I)=\sum_{j=1}^{\ell}6^{\ell-j}s_j C(I)_j.
 \tag{I.36}
\]
Call a bit $j$ coherent if $I\mapsto C(I)_j$ satisfies the
Pl\"ucker exchange equalities in I-23.  Call a seed coherent
if the diagonal action $p_I\mapsto t^{\rho_s(I)}p_I$
preserves every decomposable $r$-vector.
Then
\[
 s\text{ is coherent}
 \quad\Longleftrightarrow\quad
 s_j=0\text{ for every noncoherent bit }j.
 \tag{I.37}
\]
Thus the set of coherent seeds is a coordinate face of the
Boolean seed cube, rather than a collection of cancellations
between noncoherent code positions.

For any selected family $\mathcal S$ of coherent seeds, let
$T\subseteq[w]$ consist of the ground elements tested by its
active nonconstant bits.  By the binary classification of I-23,
each such bit is membership of one element or its complement,
so $|T|\le\ell$.  The entire family of weights depends on an
$r$-subset $I$ only through $I\cap T$.  In particular, if
$\ell\le w-2$, there are distinct $u,v\notin T$ such that
\[
 \rho_s(K\cup\{u\})=\rho_s(K\cup\{v\})
 \quad\text{for all }s\in\mathcal S
 \tag{I.38}
\]
whenever $|K|=r-1$ and $K\cap\{u,v\}=\varnothing$.
Consequently a short actual code cannot be repaired merely
by discarding its noncoherent seeds and using the surviving
weights as ground-variable cocharacters.

\paragraph{Proof of the exact filtering statement.}
For each Pl\"ucker exchange equality, write its deficit for
bit $j$ as
\[
 d_j=C(Kab)_j+C(Kcd)_j-C(Kac)_j-C(Kbd)_j.
\]
Each $d_j$ is an integer in $[-2,2]$.  Coherence of the seed
requires $\sum_j 6^{\ell-j}s_jd_j=0$ for this exchange,
and likewise for the other pairing.  If there were a largest
active power $6^h$ with nonzero coefficient, its absolute
contribution would be at least $6^h$, whereas all lower
powers together contribute at most
\[
 2\sum_{a=0}^{h-1}6^a
   =\frac25(6^h-1)<6^h.
\]
This is impossible.  Hence $s_jd_j=0$ for every $j$ and
every exchange.  An active bit must therefore be coherent.
Conversely a sum of coherent weight functions satisfies all
exchange equalities and hence preserves decomposability by
I-23.  This proves (I.37).  The common positive layer
multiplier in Hoza's formula does not affect the argument.
The deficit argument is an equality of integer exponents,
so it is valid in every coefficient-field characteristic.

For the family statement, each active coherent nonconstant
bit is a membership or complementary membership function of
one element, by I-23.  All other active bits are constant.
It follows directly that $\rho_s(I)$ is determined by
$I\cap T$.  Two untested elements exist when $|T|\le w-2$,
and exchanging them proves (I.38).  A suitable $K$ exists
since $r-1\le w-2$.

\paragraph{An explicit nonzero cancellation witness.}
Write the induced coherent weights as
$\rho_s(I)=c_s+\sum_{i\in I}a_{s,i}$.
Equation (I.38) forces $a_{s,u}=a_{s,v}$ for every seed;
this equality is independent of the common gauge ambiguity
in the choice of the $a_{s,i}$.  The nonzero polynomial
\[
 f(X_1,\ldots,X_w)=X_u-X_v
\]
therefore vanishes under every induced monomial substitution
$X_i\leftarrow Z^{a_{s,i}}$.  Adding the same scalar center
to every coordinate does not change that conclusion.
The assertion also holds for Laurent exponents; adding a
common exponent to make all exponents nonnegative preserves
the equality of the two relevant coordinates.

This is already an invertible-factor width-two example:
\[
 f=e_1^T(I_2+X_uE_{12})(I_2-X_vE_{12})e_2,
\]
with identity factors at the other positions.  Equivalently,
the decomposable vector
$e_K\wedge(X_u e_u+X_v e_v)$ has two Pl\"ucker coordinates,
and their oriented difference gives this same polynomial.
Thus the witness belongs to the algebraic setting that the
coherent decoder was intended to preserve, rather than merely
showing a combinatorial collision between unused codewords.
\hfill$\square$

\paragraph{Magnitude and scope.}
To distinguish all $r$-subsets by such a coherent family,
at least $w-1$ ground elements must be tested.  Consequently
at least $w-1$ distinct nonconstant bit positions must occur
in its seeds.  Among their powers of six, one is at least
$6^{w-2}$.  A seed activating that bit, and a codeword on
which the bit is one, give a raw weight at least $6^{w-2}$.
This last observation concerns the unmodified weights;
rescaling an individual cocharacter can change its magnitude,
so it is not asserted as an invariant degree lower bound
under arbitrary further decodings.  The exact face
description (I.37) and the short-code collision (I.38) remain
valid even if common factors of individual seeds are removed.

Longer codes followed by a separately justified small family,
or maps that leave the Grassmann cone, are not excluded.
Nor does this obstruct adding a dense SV direction: for
distinct nodes, $\lambda_u(Y)-\lambda_v(Y)$ is nonzero and
already destroys this particular linear collision.  I-26
tests that next dense repair by a different, determinant
counterexample, and leaves logarithmically many pairs open
for the general one-hidden-order rank-test class.
